% Options for packages loaded elsewhere
\PassOptionsToPackage{unicode}{hyperref}
\PassOptionsToPackage{hyphens}{url}
\PassOptionsToPackage{dvipsnames,svgnames,x11names}{xcolor}
\documentclass[
  11pt,
]{article}
\usepackage{xcolor}
\usepackage[margin=1in]{geometry}
\usepackage{amsmath,amssymb}
\setcounter{secnumdepth}{-\maxdimen} % remove section numbering
\usepackage{iftex}
\ifPDFTeX
  \usepackage[T1]{fontenc}
  \usepackage[utf8]{inputenc}
  \usepackage{textcomp} % provide euro and other symbols
\else % if luatex or xetex
  \usepackage{unicode-math} % this also loads fontspec
  \defaultfontfeatures{Scale=MatchLowercase}
  \defaultfontfeatures[\rmfamily]{Ligatures=TeX,Scale=1}
\fi
\usepackage{lmodern}
\ifPDFTeX\else
  % xetex/luatex font selection
\fi
% Use upquote if available, for straight quotes in verbatim environments
\IfFileExists{upquote.sty}{\usepackage{upquote}}{}
\IfFileExists{microtype.sty}{% use microtype if available
  \usepackage[]{microtype}
  \UseMicrotypeSet[protrusion]{basicmath} % disable protrusion for tt fonts
}{}
\makeatletter
\@ifundefined{KOMAClassName}{% if non-KOMA class
  \IfFileExists{parskip.sty}{%
    \usepackage{parskip}
  }{% else
    \setlength{\parindent}{0pt}
    \setlength{\parskip}{6pt plus 2pt minus 1pt}}
}{% if KOMA class
  \KOMAoptions{parskip=half}}
\makeatother
\usepackage{longtable,booktabs,array}
\usepackage{calc} % for calculating minipage widths
% Correct order of tables after \paragraph or \subparagraph
\usepackage{etoolbox}
\makeatletter
\patchcmd\longtable{\par}{\if@noskipsec\mbox{}\fi\par}{}{}
\makeatother
% Allow footnotes in longtable head/foot
\IfFileExists{footnotehyper.sty}{\usepackage{footnotehyper}}{\usepackage{footnote}}
\makesavenoteenv{longtable}
\setlength{\emergencystretch}{3em} % prevent overfull lines
\providecommand{\tightlist}{%
  \setlength{\itemsep}{0pt}\setlength{\parskip}{0pt}}
\usepackage{bookmark}
\IfFileExists{xurl.sty}{\usepackage{xurl}}{} % add URL line breaks if available
\urlstyle{same}
\hypersetup{
  pdftitle={Skill Loom: Evidence-Bounded Maintenance of Agent Skill Libraries},
  pdfauthor={Skill Loom contributors},
  colorlinks=true,
  linkcolor={Maroon},
  filecolor={Maroon},
  citecolor={Blue},
  urlcolor={Blue},
  pdfcreator={LaTeX via pandoc}}

\title{Skill Loom: Evidence-Bounded Maintenance of Agent Skill
Libraries}
\author{Skill Loom contributors}
\date{6 October 2026}

\begin{document}
\maketitle

\textbf{Systems technical report · version 0.1 · not peer reviewed}

\section{Abstract}\label{abstract}

Agent skills package reusable instructions, scripts and resources, but a
growing library introduces a maintenance problem distinct from
generating a useful answer. Candidate quality, user fit, host
compatibility, deployment integrity and downstream utility require
different evidence. We describe Skill Loom, a local toolkit that
separates these concerns through explicit user profiles, bounded source
discovery, capability-gap analysis, reviewed file-change plans,
recoverable journals and task-bound evidence receipts. The design is
grounded in a focused review of fifteen primary papers and official
engineering sources; it integrates established ideas rather than
claiming a new learning algorithm. We evaluate twelve deliberately
constructed software-contract scenarios and a real maintenance cycle in
one Windows Codex environment containing 39 user-maintained skills. The
reviewed transaction implementation satisfies all eight selected
mutation contracts; bound receipts satisfy all four selected evidence
contracts. These results establish only the tested software properties.
They do not demonstrate higher agent task accuracy, reduced token
expenditure, general security, or universal compatibility. We provide
reproducible code, failure cases, an anonymized skill catalog and a
separate plan for future paired model-task evaluations.

\textbf{Keywords:} agent skills; lifecycle management; provenance;
evaluation; reversible deployment; human-agent collaboration.

\section{1. Introduction}\label{introduction}

File-based agent skills expose reusable procedural knowledge through
metadata, instructions and optional resources. Progressive disclosure
allows a host to inspect descriptions before loading relevant content
{[}1{]}. However, packaging and discoverability do not establish
usefulness. SkillsBench evaluates task performance under controlled
skill conditions and reports heterogeneous effects {[}2{]}. A lifecycle
survey further distinguishes proposal, admission, retrieval, maintenance
and governance as separate concerns {[}3{]}.

The practical problem addressed here is maintaining a personal skill
library under changing tasks, users, models and harnesses. A candidate
can be relevant but unlicensed for the intended redistribution,
syntactically valid but behaviorally ineffective, or useful yet deployed
incorrectly. A rollback can restore bytes while losing executable
permissions. A receipt can reference unchanged evidence from the wrong
task. These are distinguishable failure modes, and a single quality
score cannot resolve all of them.

This report asks two bounded questions. \textbf{RQ1:} can explicit
change plans and recoverable journals preserve specified mutation and
rollback contracts under selected faults? \textbf{RQ2:} can task,
candidate and time binding reject selected stale or replayed evidence
receipts? A third question---whether the full workflow improves agent
task quality and cost across users and harnesses---remains open. The
contribution is a reproducible engineering artifact, its operational
contracts, a focused synthesis of prior work and a small fault-injection
study. We do not claim priority for the lifecycle, reflection,
deterministic-host or rollback concepts.

\section{2. Related work and design
implications}\label{related-work-and-design-implications}

\subsection{2.1 Experience, reflection and candidate
search}\label{experience-reflection-and-candidate-search}

Voyager stores reusable executable skills and refines programs using
environmental feedback {[}4{]}. Reflexion uses verbal feedback as
persistent episodic information, while Self-Refine iterates an output
through feedback and revision {[}5{]}, {[}6{]}. These mechanisms
motivate reusable diagnosis, but they do not imply that every reflection
should become a global instruction.

GEPA searches prompt candidates with reflection and validation-informed
selection {[}7{]}. EvoSkill applies a proposer/builder organization to
skill discovery {[}8{]}. SkillRL jointly changes policies and skill
libraries through training {[}9{]}. AutoSkill addresses experience-based
skill extraction and personalized library maintenance {[}10{]}. ACE
emphasizes incremental context updates and structured curation {[}11{]}.
Skill Loom borrows the separation between proposal and selection; its
manually graded ranking queue is neither a reproduction of these
optimizers nor a trained utility predictor.

\subsection{2.2 Verification, compression and negative
evidence}\label{verification-compression-and-negative-evidence}

CoEvoSkills uses co-evolving verification while retaining oracle
feedback within its reported procedure {[}12{]}. SkillZip already
describes structural preservation, provenance and deterministic
host-managed writes {[}13{]}. These are direct precedents for important
aspects of our engineering design. In both cases, verification
assumptions matter: a surrogate or a parsed structural contract is not
unrestricted semantic correctness.

Rethinking Self-Evolving Agent Skills studies multi-round feedback and
reports model- and task-dependent returns rather than guaranteed
monotonic improvement {[}14{]}. SkillCoach separates skill-use process
evaluation from final outcomes {[}15{]}. A broader evaluation survey
organizes mechanisms and benchmark coverage {[}16{]}. Together, these
works support preserving rejected candidates, no-change outcomes and
honest stopping states. They also argue against using installation count
or self-reported success as the optimization target.

\subsection{2.3 Operational maintenance}\label{operational-maintenance}

The technical-debt perspective highlights hidden dependencies and
feedback loops in maintained learning systems {[}17{]}. We use this as
an engineering analogy, not as evidence about SKILL.md performance.
Official agent-evaluation guidance distinguishes tasks, trials,
transcripts, outcomes and graders {[}18{]}. Context-engineering guidance
motivates selective retrieval and clear tool responsibilities {[}19{]}.
Accordingly, Skill Loom separates static inspection, host discovery,
script checks and task-level evaluation.

The accompanying literature matrix fixes paper versions and gives
section-level evidence locations. It notes incompatible benchmark
versions, limited replication and reported inconsistencies instead of
constructing a cross-paper leaderboard. None of the cited model
experiments was reproduced in this study.

\section{3. Problem formulation}\label{problem-formulation}

Let an execution setting be \(E=(U,M,H,T,P,D,V,B)\), comprising a user
profile, model, harness, tools, permissions, task distribution,
evaluator and budget. A library state \(L_t\) includes skill content,
source identities, local adaptations and deployment metadata. A proposal
mechanism produces a candidate \(C_t\) from selected task evidence.
Admission either retains the existing library or applies a reviewed
transition:

\[L_{t+1}=A(L_t,C_t,E).\]

Retention is a legitimate result. Three properties must be
distinguished: \textbf{proposal relevance}, \textbf{state-transition
correctness}, and \textbf{behavioral utility}. A verified file-state
transition does not establish proposal relevance or behavioral utility,
and a useful candidate does not excuse an incorrect state transition.

For a task set \(D\), one may seek a portfolio balancing quality \(Q\),
cost \(C\) and maintenance burden \(K\):

\[\max_L\; Q(L;E,D)-\lambda C(L;E,D)-\mu K(L),\]

subject to user authorization, output requirements and compatibility
constraints. This expression defines a design objective, not an
implemented optimizer or a proven convergence result. The prototype's
profile weights order an investigation queue; unknown evidence remains
unknown and hard adoption conditions are reported separately.

\subsection{3.1 Transition contracts and
assumptions}\label{transition-contracts-and-assumptions}

A reviewed plan contains before/after file hashes, supported permission
modes, exact changes, a physical target root and a disjoint journal
location. Application rechecks both snapshots, including a no-change
plan. Modified originals are backed up before writes. A completed
transition must match its after state. An interrupted transition retains
a journal for recovery. Rollback refuses later user changes rather than
overwriting them.

Under cooperative exclusive writing, available storage and successful
filesystem operations, these checks provide a direct argument for the
tested byte-and-mode postconditions: each planned path receives the
reviewed content and supported mode, and the final snapshot is compared
with the plan. This is not a proof of crash-consistent multi-file
atomicity. ACLs, extended attributes, directory metadata, adversarial
same-user races and arbitrary storage failures are outside the claim.
Individual file replacement is used; a multi-file transition may be
interrupted and require explicit recovery.

\section{4. System design}\label{system-design}

\subsection{4.1 User needs and bounded
discovery}\label{user-needs-and-bounded-discovery}

An editable private profile records required capabilities and
user-selected budgets. A lightweight environment check reports OS,
Python and executable presence; it does not scan credentials or infer
preferences. Capability observations distinguish skill, environment,
source and unknown causes. Existing providers are considered before
searching or creating a skill. An unobserved skill is not automatically
removed.

Source discovery observes registered repositories and optionally
searches explicit public keywords. Coverage is the registered source set
and bounded search sample, not the entire internet. Staging requires a
full commit, verifies downloaded Git blob identities and lengths,
retains source hashes and license files, and does not execute downloaded
code. Five evidence-backed dimensions support candidate review: task
fit, evidence, maintainability, efficiency and portability. These
dimensions are configurable policy, not calibrated probabilities.

\subsection{4.2 Portfolio and
deployment}\label{portfolio-and-deployment}

The catalog maps skills to capability identifiers and scope. Capability
overlap and user-defined metadata budgets identify review candidates.
Similarity does not establish redundancy: two skills may serve different
artifacts or authorization boundaries. Merging, narrowing, retiring or
keeping an item requires examining representative tasks and its
consumers.

Deployment is a separate command from discovery and ranking. The plan
digest selects the reviewed plan; it is an integrity identifier, not a
signature or grant of authority. Journals reside outside the discovery
root. Discovery can inspect a linked entry, whereas mutation requires
the reviewed physical skills root. The prototype does not change host
approval settings or silently install hooks.

\subsection{4.3 Evidence receipts and bounded
repair}\label{evidence-receipts-and-bounded-repair}

Receipts enumerate required checks and reference evidence bytes.
Optional CLI constraints bind a receipt to a run identifier, candidate
digest and age limit; the supplied Stop-hook example requires run and
candidate bindings. A failed first Stop may request one scoped repair. A
repeated Stop or an unavailable gate permits termination while
communicating an unverified state. Allowing termination never converts
failure to success.

The receipt checker verifies declarations and artifacts, not their
truthfulness. A fabricated report with matching hashes remains a
fabricated report. Grader correctness, real command execution and the
independence of an expected candidate identifier require external
controls and review.

\subsection{4.4 Portability and data
management}\label{portability-and-data-management}

The common layer is file-based and independent of any model API.
Host-specific discovery paths, hook semantics and delegation remain
explicit adapters. Missing subagent support falls back to serial work;
unavailable hooks fall back to explicit gate execution. Protocol
compatibility and behavioral transfer are separate evaluation targets.

Raw traces, private profiles and journals remain outside the public
repository. The optional trace observer emits structural counts without
prompts, commands or tool outputs. Cache cleanup accepts only explicitly
adopted reconstructible directories, rechecks exact files and records
deletion intent before unlinking. It does not manage credentials,
sessions, databases or plugin caches.

\section{5. Evaluation protocol}\label{evaluation-protocol}

\subsection{5.1 Deterministic fault
injection}\label{deterministic-fault-injection}

The study uses twelve hand-selected scenarios, each in a fresh temporary
directory. Eight mutation scenarios cover a clean update, missing
reference, Python syntax error, candidate drift, runtime drift, an added
local file, post-install rollback drift and a corrupt backup. Four
receipt scenarios cover a valid record, changed artifact, another task's
receipt and an expired receipt.

The comparators are small reference implementations: unchecked copying,
static inspection followed by copying, and declared-pass receipt
acceptance. They are intentionally weak operational baselines, not
published skill-evolution systems or production package managers. The
tested contract determines whether application/acceptance or refusal is
expected. Source code and per-case outcomes are provided in
\texttt{scripts/benchmark\_invariants.py} and
\texttt{research/results/invariants.json}.

Scenarios are the descriptive units. Repeated deterministic executions
are not independent samples; no confidence interval, significance test
or population success probability is estimated. The cases were chosen
during development and therefore do not constitute an independent
held-out benchmark. The study tests mechanism behavior, not
generalization.

\subsection{5.2 Local deployment study}\label{local-deployment-study}

A single Windows environment with Python 3.12.8 and Codex CLI 0.160.1
contained 39 user-maintained skills. We inventoried them and applied two
reviewed file changes to its existing maintenance skill, adding an
on-demand reference to the new workflow. We then restored the original
state, compared it with the pre-change snapshot, reapplied the candidate
and checked a no-change run. Global host configuration and
instruction-file hashes were compared before and after. This is a case
study of deployment mechanics, not a representative sample of users.

\section{6. Results}\label{results}

\begin{longtable}[]{@{}
  >{\raggedright\arraybackslash}p{(\linewidth - 4\tabcolsep) * \real{0.3000}}
  >{\raggedright\arraybackslash}p{(\linewidth - 4\tabcolsep) * \real{0.3000}}
  >{\raggedleft\arraybackslash}p{(\linewidth - 4\tabcolsep) * \real{0.4000}}@{}}
\toprule\noalign{}
\begin{minipage}[b]{\linewidth}\raggedright
Contract family
\end{minipage} & \begin{minipage}[b]{\linewidth}\raggedright
Reference implementation
\end{minipage} & \begin{minipage}[b]{\linewidth}\raggedleft
Satisfied / selected cases
\end{minipage} \\
\midrule\noalign{}
\endhead
\bottomrule\noalign{}
\endlastfoot
File mutation and rollback & Unchecked copy & 1 / 8 \\
File mutation and rollback & Static check plus copy & 3 / 8 \\
File mutation and rollback & Reviewed transaction & 8 / 8 \\
Receipt acceptance & Declared pass only & 1 / 4 \\
Receipt acceptance & Task/candidate/time-bound receipt & 4 / 4 \\
\end{longtable}

The static comparator rejects missing references and syntax errors but
does not enforce reviewed-state or rollback conditions. The transaction
implementation satisfies the eight selected contracts. Bound receipts
reject the three selected stale/replay cases while accepting the valid
record. These differences are expected consequences of the implemented
checks and should not be interpreted as broad empirical superiority.

In the local deployment cycle, the skill count remained 39; the final
static inventory reported zero errors and zero warnings. Rollback
matched the original file snapshot, reapplication completed, and the
no-change run returned without changing content. The protected host
configuration hashes remained unchanged. No new same-purpose maintenance
skill was added to the global discovery catalog.

Independent review exposed additional implementation gaps: executable
permission loss on POSIX, missing drift checks on no-change plans,
incomplete cleanup interruption records, weak input-type checks, missing
Git blob verification, nonfinite freshness limits and a case-insensitive
reserved-name collision. The implementation and regression tests were
amended. The latest software-test and CI status is separately recorded
in \texttt{docs/validation.md}; platform-skipped tests are not counted
as executed checks. Native Claude Code hook operation and broad
model-task transfer remain unverified in this local study.

\section{7. Limitations and next
experiments}\label{limitations-and-next-experiments}

The project does not establish an automatic quality oracle, worldwide
discovery coverage or downstream model improvement. Its source ranking
depends on reviewed annotations. Capability maps can be incomplete.
Read-only syntax checks do not detect arbitrary malicious code. Hashes
do not validate semantics. The filesystem protocol assumes cooperative
writers and does not preserve every OS-specific metadata field. The
scope is a local prototype rather than a multi-tenant managed service.

A subsequent behavioral study should compare no skill, the current
stable skill and the candidate on matched tasks, fixing model, harness,
tools and budget. Selection tasks must be separate from frozen
evaluation tasks. Process measures should include correct and spurious
skill selection, while outcome measures should include artifact
fidelity, success, retries, latency and total cost. A multi-round study
should count rejected candidates, no-ops and retained old versions, and
compare equal-budget alternatives. Model or user migration requires new
tests rather than inherited confidence.

These experiments are proposed, not completed. Larger sample sizes
should follow an explicit research question, a variance estimate and a
feasible budget; this report does not invent expected improvements or a
power calculation.

\section{8. Conclusion}\label{conclusion}

Skill Loom makes a limited but useful distinction operational:
generating a candidate, changing a library correctly and improving a
user's tasks are different achievements. Its local toolkit supports
inspectable decisions, bounded repair and reversible deployment. The
reported fault-injection and local case-study results support specified
software contracts, while leaving model utility, broad portability and
comparative usability as explicit future evaluation questions.

\section{Availability and disclosure}\label{availability-and-disclosure}

Code, anonymized examples, source records, experiment scripts and this
manuscript are released with the project. Third-party skill content is
not redistributed by the catalog. Development used an AI coding
assistant and independent agent reviews; the promotional card is
AI-generated. The technical report has not undergone academic peer
review and does not imply acceptance by a venue. No private user
transcripts or credentials are included in the public artifact.

\section{References}\label{references}

{[}1{]} B. Zhang, K. Lazuka, and M. Murag. Equipping agents for the real
world with Agent Skills. Anthropic Engineering, 2025.
\url{https://www.anthropic.com/engineering/equipping-agents-for-the-real-world-with-agent-skills}
{[}2{]} X. Li et al.~SkillsBench: Benchmarking How Well Agent Skills
Work Across Diverse Tasks. arXiv:2602.12670v4, 2026.
\url{https://arxiv.org/abs/2602.12670v4} {[}3{]} Y. Li. Dynamic Agent
Skills: A Lifecycle Survey and Taxonomy of Evolving Skill Libraries.
arXiv:2607.10113v1, 2026. \url{https://arxiv.org/abs/2607.10113v1}
{[}4{]} G. Wang et al.~Voyager: An Open-Ended Embodied Agent with Large
Language Models. arXiv:2305.16291v2, 2023.
\url{https://arxiv.org/abs/2305.16291v2} {[}5{]} N. Shinn et
al.~Reflexion: Language Agents with Verbal Reinforcement Learning.
arXiv:2303.11366v4, 2023. \url{https://arxiv.org/abs/2303.11366v4}
{[}6{]} A. Madaan et al.~Self-Refine: Iterative Refinement with
Self-Feedback. arXiv:2303.17651v2, 2023.
\url{https://arxiv.org/abs/2303.17651v2} {[}7{]} L. A. Agrawal et
al.~GEPA: Reflective Prompt Evolution Can Outperform Reinforcement
Learning. arXiv:2507.19457v2, 2025.
\url{https://arxiv.org/abs/2507.19457v2} {[}8{]} S. Alzubi et
al.~EvoSkill: Automated Skill Discovery for Multi-Agent Systems.
arXiv:2603.02766v1, 2026. \url{https://arxiv.org/abs/2603.02766v1}
{[}9{]} P. Xia et al.~SkillRL: Evolving Agents via Recursive
Skill-Augmented Reinforcement Learning. arXiv:2602.08234v1, 2026.
\url{https://arxiv.org/abs/2602.08234v1} {[}10{]} Y. Yang et
al.~AutoSkill: Experience-Driven Lifelong Learning via Skill
Self-Evolution. arXiv:2603.01145v2, 2026.
\url{https://arxiv.org/abs/2603.01145v2} {[}11{]} Q. Zhang et
al.~Agentic Context Engineering: Evolving Contexts for Self-Improving
Language Models. arXiv:2510.04618v3, 2025.
\url{https://arxiv.org/abs/2510.04618v3} {[}12{]} H. Zhang et
al.~CoEvoSkills: Self-Evolving Agent Skills via Co-Evolutionary
Verification. arXiv:2604.01687v3, 2026.
\url{https://arxiv.org/abs/2604.01687v3} {[}13{]} X. Bai et
al.~SkillZip: Evaluation-Free Skill Compression for Self-Evolving Agents
by Discovering Reusable Structure. arXiv:2608.11079v2, 2026.
\url{https://arxiv.org/abs/2608.11079v2} {[}14{]} Y. Liu et
al.~Rethinking Self-Evolving Agent Skills: Feedback Dynamics over
Multiple Rounds. arXiv:2608.02636v1, 2026.
\url{https://arxiv.org/abs/2608.02636v1} {[}15{]} J. Zhu et
al.~SkillCoach: Self-Evolving Rubrics for Evaluating and Enhancing
Agentic Skill-Use. arXiv:2607.01874v1, 2026.
\url{https://arxiv.org/abs/2607.01874v1} {[}16{]} K. Ding et al.~Agent
Skill Evaluation and Evolution: Frameworks and Benchmarks.
arXiv:2606.11435v1, 2026. \url{https://arxiv.org/abs/2606.11435v1}
{[}17{]} D. Sculley et al.~Hidden Technical Debt in Machine Learning
Systems. NeurIPS, 2015, pp.~2503--2511.
\url{https://papers.neurips.cc/paper/5656-hidden-technical-debt-in-machine-learning-systems.pdf}
{[}18{]} Anthropic. Demystifying evals for AI agents. Anthropic
Engineering, 2026.
\url{https://www.anthropic.com/engineering/demystifying-evals-for-ai-agents}
{[}19{]} Anthropic. Effective context engineering for AI agents.
Anthropic Engineering, 2025.
\url{https://www.anthropic.com/engineering/effective-context-engineering-for-ai-agents}

\end{document}
